% LTeX: language=de-DE
\chapter{Kurzfassung}

Legacy-Softwareentwicklungsumgebungen stellen erhebliche Herausforderungen dar, die die Produktivität der Entwickler beeinträchtigen und die Reproduzierbarkeit der Builds gefährden.
Herkömmliche Entwicklungsumgebungen basieren auf manuell konfigurierten Workstations, was zu Konfigurationsabweichungen und dem berüchtigten Problem „Auf meinem Rechner funktioniert es“ führt.
Neue Teammitglieder müssen unter Umständen tagelang ihre Umgebung konfigurieren, bevor sie produktiv arbeiten können, während virtuelle Maschinen einen erheblichen Ressourcenaufwand verursachen.

Diese Arbeit untersucht die Anwendung der Containerisierung auf die Entwicklung von Legacy-Software und untersucht, wie Docker-basierte Containertechnologien die Entwicklungsabläufe für langjährige native Softwareprojekte modernisieren können.
Angewandt auf ein industrielles C++-Softwaresystem mit jahrzehntelang gewachsener Komplexität liefert diese Forschung empirische Belege für den Vergleich von containerisierten Entwicklungsumgebungen mit traditionellen nativen und virtualisierten Ansätzen hinsichtlich Build-Zeiten, Ressourcenauslastung und Entwicklerproduktivitätsmetriken.

\begin{sloppypar}
Die Ergebnisse zeigen, dass sich Container durch automatisierte Bereitstellung mithilfe von Infrastructure-as-Code, nahtlose IDE-Integration über DevContainers und standardisierte CMake-basierte plattformübergreifende Build-Konfigurationen in Legacy-Workflows integrieren lassen.
Die Umstellungszeit für Toolchains wurde um 95\,\% reduziert---von 12,9 Minuten auf 37,5 Sekunden---, was direkte Produktivitätssteigerungen für Entwickler-Onboarding und Multi-Version-Entwicklungsszenarien bestätigt.
\end{sloppypar}

\begin{sloppypar}
Die Leistungsbewertung ergab, dass Linux-basierte Container bei dateisystemintensiven Vorgängen um 35 bis 72\,\% besser abschnitten als Windows-native Ansätze, während bei CPU-gebundener Kompilierung in allen Umgebungen nur geringe Unterschiede von maximal 15\,\% zu verzeichnen waren.
Native Linux-Container verursachten im Vergleich zur Bare-Metal-Ausführung nur einen vernachlässigbaren Overhead.
Das Compiler-Caching verstärkte die Leistungsunterschiede erheblich, wobei optimierte containerisierte Workflows im Vergleich zu herkömmlichen, nicht optimierten Ansätzen eine Reduzierung der Gesamt-Workflow-Zeit um bis zu 89\,\% erzielten.
Die Analyse der Ressourcennutzung zeigte, dass die CPU-Auslastungsmuster eher die Art der Operation als die Art der Umgebung widerspiegelten, während der Speicherverbrauch in Hypervisor-basierten Lösungen im Vergleich zur nativen Containerisierung 5 bis 10 GB zusätzliche Ressourcen erforderte.
\end{sloppypar}

Die Resultate belegen, dass sich die Prinzipien der Containerisierung erfolgreich von der Bereitstellung von Microservices auf monolithische Legacy-Systeme übertragen lassen.
Dadurch werden Reproduzierbarkeit, Produktivität und Performance deutlich verbessert.
Containerisierung stellt somit eine praktikable Modernisierungsstrategie dar, um bestehende Software-Investitionen zu erhalten und gleichzeitig wettbewerbsfähige Entwicklungspraktiken sicherzustellen.